%%%%%%%%%%%%%%%%%%%%%%%%%%%%%%%%%%%%%%%%%
% Medium Length Graduate Curriculum Vitae
% LaTeX Template
% Version 1.1 (9/12/12)
%
% This template has been downloaded from:
% http://www.LaTeXTemplates.com
%
% Original author:
% Rensselaer Polytechnic Institute (http://www.rpi.edu/dept/arc/training/latex/resumes/)
%
% Important note:
% This template requires the res.cls file to be in the same directory as the
% .tex file. The res.cls file provides the resume style used for structuring the
% document.
%
%%%%%%%%%%%%%%%%%%%%%%%%%%%%%%%%%%%%%%%%%

%----------------------------------------------------------------------------------------
%	PACKAGES AND OTHER DOCUMENT CONFIGURATIONS
%----------------------------------------------------------------------------------------

\documentclass[margin, 9.5pt]{res} % Use the res.cls style, the font size can be changed to 11pt or 12pt here

\usepackage{helvet} % Default font is the helvetica postscript font
% \usepackage{newcent} % To change the default font to the new century schoolbook postscript font uncomment this line and comment the one above

\setlength{\textwidth}{5.3in} % Text width of the document
\usepackage{xcolor}
\usepackage{hyperref}
\usepackage{array}
\hypersetup{
    colorlinks=true,
    linkcolor={black!45!blue},
    urlcolor={black!45!blue},
}

\begin{document}

%----------------------------------------------------------------------------------------
%	NAME AND ADDRESS SECTION
%----------------------------------------------------------------------------------------

\moveleft.5\hoffset\centerline{\huge\bf Tom Bearpark} % Your name at the top
 
\moveleft\hoffset\vbox{\hrule width\resumewidth height 1pt}\smallskip % Horizontal line after name; adjust line thickness by changing the '1pt'
 
\moveleft.5\hoffset\centerline{\href{mailto:t.bearpark@gmail.com}{t.bearpark@gmail.com} \quad  \quad \href{https://tombearpark.com/}{tombearpark.com}}

%----------------------------------------------------------------------------------------

\begin{resume}

\section{CURRENT\\POSITIONS}

{\sl Post-doctoral Research Fellow} \hfill 2026-present\\
University of Exeter \\
Land, Environment, Economics and Policy Institute (LEEP).

{\sl Visiting Fellow} \hfill 2026-present\\
London School of Economics and Political Science\\
Grantham Research Institute on Climate Change and the Environment.
 

\moveleft\hoffset\vbox{\color{gray}\hrule width\resumewidth height .5pt}\par\vspace{-1em}
\section{EDUCATION}

{\sl Ph.D, Public Affairs} \hfill  2025\\
Princeton University\\
{\small Fields: Environmental economics, Climate impacts, Econometrics, Quantitative methods.} 

{\sl MPhil, Economic Research} \hfill 2018\\
University of Cambridge

% {\sl Post Graduate Certificate, Regulation of the Power Sector} \hfill 2016\\
% European University Institute

{\sl BSc, Economics and Philosophy} \hfill 2015\\
London School of Economics and Political Science

\moveleft\hoffset\vbox{\color{gray}\hrule width\resumewidth height .5pt}\par\vspace{-1em}
\section{PEER REVIEWED WORK}

\textbf{Causes and Correlates of Unsustainable Scuba Diving Tourism on Coral Reefs} \\
2026, \textit{Conservation Letters}. \\
\textit{
    (with 
    \href{https://www.thebinglin.com/}{Bing Lin}, 
    \href{https://www.bangor.ac.uk/staff/sos/ronan-roche-076894/en}{Ronan C. Roche}, 
    \href{https://www.linkedin.com/in/jacobo-mcguire-14bb9b68/}{Jacobo McGuire}, 
    \href{https://rare.org/inspiring-human-nature/raymond-jakub/}{Raymond Jakub}, 
    \href{https://www.researchgate.net/profile/Derta-Prabuning}{Derta Prabuning}, 
    \href{https://elke-u-weber.com/}{Elke U. Weber}, and \href{https://eeb.princeton.edu/people/david-wilcove}{David S. Wilcove})
} \\
\url{https://doi.org/10.1111/con4.70055}


\textbf{Mortality impacts of rainfall and sea level rise in a developing megacity} \\
2025, \textit{Nature}. \\
\textit{
    (with 
    \href{https://sites.google.com/site/ashwinrode/home}{Ashwin Rode} and 
    \href{https://scholar.google.co.in/citations?user=KOWEt48AAAAJ&hl=en}{Archana Patankar})
} \\
\url{https://doi.org/10.1038/s41586-025-09730-4}

\textbf{Data Anomalies and the Economic Commitment of Climate Change}\\
2025, \textit{Nature}. \\
\textit{
    (with 
    \href{https://dylanhogan.github.io/}{Dylan Hogan}, and 
    \href{https://www.solomonhsiang.com/}{Solomon Hsiang})
}\\
\url{https://doi.org/10.1038/s41586-025-09320-4}

\moveleft\hoffset\vbox{\color{gray}\hrule width\resumewidth height .5pt}\par\vspace{-1em}

\section{WORKING PAPERS}

\textbf{Modelling climate migration} \\
\textit{
    (with 
    \href{https://nchoquettelevy.github.io/}{Nic Choquette-Levy}, 
    \href{https://oppenheimer.scholar.princeton.edu/}{Michael Oppenheimer}, 
    \href{https://jrosenthalkay.github.io/}{Jordan Rosenthal-Kay}, and 
    \href{https://cpree.princeton.edu/people/tingyin-xiao}{Tingyin Xiao})
}\\
% Preparing for submission to \textbf{\textit{Nature}}. \\
Draft available \href{https://www.dropbox.com/scl/fi/350yvqv8rzqh9vi14s8hw/DRAFT.pdf?rlkey=d1ru0tq0fogmaw3rccr5uyw3t&st=x0mxyoea&dl=0}{here}. 

\small
% \textit{Abstract}: Quantifying and projecting climate migration is a critical challenge for policymakers and researchers. While causal inference, agent-based, gravity, and general equilibrium models have been employed widely to estimate the climate-migration relationship, their results often diverge, and they are typically used in isolation and at different spatio-temporal scales. Here, we aim to bridge this gap by introducing a unified notation and conceptual framework, designed to compare and evaluate the assumptions and outputs of these diverse models at a common scale. We illustrate our comparison using an empirically relevant case study: estimating and projecting the relationship between annual average temperature and internal state-to-state migration in the United States. We analyse results from hundreds of model permutations, illustrating how methodological choices and researcher degrees of freedom shape model outcomes. Our findings show that variations in modelling approach introduce uncertainty comparable in magnitude to statistical uncertainty. By integrating policy-relevant insights across methods, we highlight the value of model intercomparison in improving projections and informing policy decisions.
\normalsize

\textbf{Selecting time controls in climate impact studies} \\
\textit{
    (with \href{https://filippopalomba.github.io/}{Filippo Palomba})
    }\\
% Preparing for submission at a general interest economics journal (AEJ Applied). \\
Draft available \href{https://www.dropbox.com/scl/fi/n1j6o8bfq6zw2b02rh6pr/MAIN.pdf?rlkey=euho0xbt3qmiiv6dgqhx3tcgx&st=g4jl2qvf&dl=0}{here}. 

\small
% \textit{Abstract}: Empirical climate impacts research is dominated by studies that use panel fixed-effects regressions to estimate the causal effect of weather on outcomes. However, the functional form of time controls in these studies, i.e., the specification of time trends, lacks a clear theoretical foundation or data-driven justification. Moreover, these choices can substantially influence the magnitude of estimated results in important applications. In this paper, we elucidate a framework for choosing time controls in climate impacts studies. We propose a two-equation model which aligns with the reasoning used in climate impacts studies to justify identification. In our framework, competing models are evaluated according to their ability to isolate plausibly random variation in the treatment variable. We provide simulations to illustrate our proposal, and apply our framework to an open and policy relevant agenda in the climate impacts literature; estimating the effect of temperature on GDP growth rates. 
\normalsize

\textbf{The economic geography of climate risk} \\
\textit{
    (with
    \href{https://www.aditya-bhandari.com/}{Aditya Bhandari} and
    \href{https://jrosenthalkay.github.io/}{Jordan Rosenthal-Kay})
} \\
Draft available \href{https://www.dropbox.com/scl/fi/tkq8wntapi4m3h527znep/EGCR.pdf?rlkey=fapv9duz9gm32ntjqfxtv4gsk&st=jpji23ts&dl=0}{here}.

\textbf{Estimating turning points and testing for climate adaptation} \\
\textit{
    (with \href{https://jrosenthalkay.github.io/}{Jordan Rosenthal-Kay})
} \\
Draft available \href{https://tombearpark.com/assets/2026_turning_points.pdf}{here}.

\textbf{The Carbon at Risk measure to manage carbon removal risk and guide effective portfolio construction} \\
\textit{
    (with
    \href{https://www.carbonatrisk.org/bali-lee-1}{Bali Lee},
    Nick Gogerty,
    Oleksandr Kit,
    Delia Meth-Cohn,
    Frederic Olbert,
    \href{https://www.lse.ac.uk/granthaminstitute/profile/frank-venmans/}{Frank Venmans},
    Gabrielle Walker,
    Paul Young, and
    \href{https://experts.exeter.ac.uk/35112-ben-groom}{Ben Groom})
} \\
Draft available \href{https://www.researchsquare.com/article/rs-9382336/v1}{here}.

\moveleft\hoffset\vbox{\color{gray}\hrule width\resumewidth height .5pt}\par\vspace{-1em}
\section{SELECTED WORK IN PROGRESS}

\textbf{The impact of climate change on global economic output} \\
\textit{
    (with 
    \href{https://web.stanford.edu/~mburke/}{Marshall Burke}, 
    \href{https://dylanhogan.github.io/}{Dylan Hogan}, and 
    \href{https://www.solomonhsiang.com/}{Solomon Hsiang})
}

\textbf{The impact of urban floods on housing markets: Evidence from Mumbai} \\ 
\textit{
    (with \href{https://aouazad.github.io/Website/}{Amine Ouazad}, 
    \href{https://sites.google.com/site/ashwinrode/home}{Ashwin Rode} and 
    \href{https://sites.google.com/view/vaidehitandel/home}{Vaidehi Tandel})
}

% \textbf{Groundwater depletion reduces farm worker employment in the High Plains} \\ 
% \textit{
%     (with 
%     \href{https://cpree.princeton.edu/people/lucas-frye}{Lucas Frye} and 
%     \href{https://oppenheimer.scholar.princeton.edu/}{Michael Oppenheimer})
% }

\moveleft\hoffset\vbox{\color{gray}\hrule width\resumewidth height .5pt}\par\vspace{-1em}
\section{FELLOWSHIPS}
{\sl 2025} - \textit{Princeton Institute for International and Regional Studies Fellowship}\\
{\sl 2025} - \textit{Princeton Graduate Fellowship in Social Data Science}\\
{\sl 2024 and 2025} - \textit{Neil A McConnel fellow}

\moveleft\hoffset\vbox{\color{gray}\hrule width\resumewidth height .5pt}\par\vspace{-1em}
\section{TEACHING}
Econometrics and Public Policy (Advanced)
\begin{itemize} \itemsep -2pt
    \item Masters level course in the School of Public and International Affairs. 
    \item Applied econometrics using R, with a focus on the concepts and techniques most relevant to public policy analysis. 
\end{itemize}

Climate Change: Impacts, Adaptation, Policy
\begin{itemize} \itemsep -2pt
    \item Undergraduate level course in the Department of Geosciences. 
    \item An exploration of the consequences of human-induced climate change and their implications for policy.
\end{itemize}
%----------------------------------------------------------------------------------------
%	PROFESSIONAL EXPERIENCE SECTION
%----------------------------------------------------------------------------------------
% \section{RESEARCH \\ EXPERIENCE}

% {\sl Research Assistant, Marianne Bertrand} \hfill 2021\\
% University of Chicago, Booth School of Business, Chicago IL
% \begin{itemize}
%     \item Conducted analysis on the political influence of institutional investors, developing and implementing machine learning algorithms to enhance the accuracy of detecting and classifying investor political activities 
% \end{itemize}

% {\sl Pre-doctoral Fellow} \hfill 2016 - 2019 \\
% Energy Policy Institute at the University of Chicago, Chicago IL
% \begin{itemize} \itemsep -2pt % Reduce space between items
% \item Assisted Professor Thomas Covert in estimating firm productivity outcomes and analyzing allocative efficiency within the oil and gas sector
% \item Worked with Professor Fiona Burlig to build panel data on economic census covariates in India to estimate the impacts of rural electrification on the manufacturing sector 
% \end{itemize}
 
% {\sl Research Experience} \hfill 2015 - 2016\\
% Columbia University, New York, NY
% \begin{itemize} \itemsep -2pt
% \item Analyzed Census data in STATA to support Professor Douglas Almond in our economics paper on sex selection
% \item Reviewed drafts and assisted in formulating new directions for existing papers
% \end{itemize} 
\moveleft\hoffset\vbox{\color{gray}\hrule width\resumewidth height .5pt}\par\vspace{-1em}
\section{WORK\\ EXPERIENCE}

{\sl Pre-doctoral fellow, \href{https://impactlab.org/}{Climate Impact Lab}} \hfill 2018 - 2020 \\
Energy Policy Institute at the University of Chicago
% \begin{itemize} \itemsep -2pt
%     \item Worked as research assistant in a large interdisciplinary team of researchers that aims to quantify the socio-economic costs of climate change. 
%     % \item Contributed to papers assessing and projecting the impacts of global climate change on migration patters, energy consumption, and mortality rates. 
%     % \item Wrote code to analyse large spatial datasets. Plotted maps and graphs to provide visualisations of results.
% \end{itemize} 

{\sl Economist, Office for Gas and Electricity Markets} \hfill 2015 - 2017 \\
UK Civil Service
% \begin{itemize} \itemsep -2pt
%     \item Researched drivers of change in the energy sector, and helped develop the `State of the Market' report, which describes the state of Great Britain's energy markets. 
%     % \item Completed a three month secondment to the Brussels office of the Council of European Energy Regulators.
% \end{itemize} 

% \textcolor{gray}{\rule{\resumewidth}{0.5pt}}\par\vspace{-\parskip}
\moveleft\hoffset\vbox{\color{gray}\hrule width\resumewidth height .5pt}\par\vspace{-1em}
\section{CONFERENCES \& INVITED SEMINARS} 
{\sl 2026:} Richard Doll Seminar (University of Oxford, UK), Energy and Environment Group Talk Series (University of Cambridge).\\
\\
{\sl 2025:} Scenarios Forum (Leeds University, UK), Managed retreat (Columbia University, USA), Eastern Economics Association (New York, USA), UK Network of Environmental Economists (London, UK), Econometric Society Interdisciplinary Frontiers (Barcelona School of Economics, Spain).\\
\\
{\sl 2024:} American Geophysical Union (Washington DC, USA), Mid-West Economics Association (Chicago, USA), Interdisciplinary Ph.D. Workshop in Sustainable Development (Columbia University, USA).
% {\sl 2023:}  

\moveleft\hoffset\vbox{\color{gray}\hrule width\resumewidth height .5pt}\par\vspace{-1em}
\section{OTHER} 

{\sl Computing:} R, Python, Julia, Stata, Git, Bash. \\[0.3em]  
% {\sl Data Analysis:} Causal inference, machine learning, spatial data, data visualisation.\\[0.3em]
{\sl Reviewing:} Climactic Change, Review of Economics and Statistics, Environment and Development Economics, Earth's Future, PNAS, UnJournal. 


\moveleft\hoffset\vbox{\color{gray}\hrule width\resumewidth height .5pt}\par\vspace{-1em}
\section{REFERENCES}

\noindent
% \small
\begin{minipage}[t]{0.47\resumewidth}
\textbf{Professor \href{https://oppenheimer.scholar.princeton.edu/}{Michael Oppenheimer}}~\\ 
  Princeton University\\
  School of Public and International Affairs\\[0.3em]

  \textbf{Professor \href{https://www.filizgarip.com/}{Filiz Garip}}\\
  Princeton University\\
  Department of Sociology
\end{minipage}%
\hfill
\begin{minipage}[t]{0.47\resumewidth}
  \textbf{Professor \href{https://www.solomonhsiang.com/}{Solomon Hsiang}}\\
  Stanford University\\
  Doerr School of Sustainability\\[0.3em]

  \textbf{Professor \href{https://www.princeton.edu/~davidlee/}{David Lee}}\\
  Princeton University\\
  Department of Economics
\end{minipage}



\end{resume}
\end{document}